
\paragraph{A}
Before formalizing the cost-minimizing choice of technique, we establish the descriptive macro context for the United States across five distinct historical windows: the pre-war Fordist recovery (1931–1941), the World War II mobilization block (1942–1945), the postwar Fordist Golden Age (1947–1973), the post-Fordist restructuring era prior to the Global Financial Crisis (1974–2007), and the post-GFC stagnation period (2008–2024). We track the average annual growth rates (mean annual log differences) of real corporate output alongside gross capital stocks built using the Gross Perpetual Inventory Method (GPIM), distinguishing private machinery and equipment (\(K^{\text{ME}}\)), private nonresidential structures (\(K^{\text{NRC}}\)), and public infrastructure in the form of government transportation, highways, and streets (\(K^{\text{GT}}\)). We report both Non-Financial Corporate sector Gross Value Added (\(y_t\)) and total economy Real GDP (\(y_{\text{total}}\)) to capture corporate-level dynamics relative to aggregate activity.

\begin{table}[H]
	\centering
	\caption{Descriptive Macro Indicators for the United States (Growth Rates, \% per year)}
	\label{tab:descriptive_macro_indicators_for_the_united_stat}
	\begin{tabular}{p{5.5cm}p{2.2cm}p{2.2cm}p{2.2cm}p{2.2cm}p{2.2cm}}
		\toprule
		Variable & Pre-war Fordism (1931–1941) & WWII Block (1942–1945) & Postwar Fordism (1947–1973) & Post-Fordism Pre-GFC (1974–2007) & Post-GFC (2008–2024) \\
		\midrule
		Real Corporate GVA (\(y_t\)) & 5.80\% & 4.36\% & 4.98\% & 3.55\% & 1.99\% \\
		Real GDP (Total Economy, \(y_{\text{total}}\)) & 3.94\% & 9.92\% & 3.76\% & 2.97\% & 1.95\% \\
		Gross Machinery \& Equipment (\(K^{\text{ME}}\)) & 1.37\% & 4.73\% & 8.49\% & 6.58\% & 3.66\% \\
		Gross Nonresidential Structures (\(K^{\text{NRC}}\)) & 1.08\% & 3.79\% & 6.38\% & 7.02\% & 3.99\% \\
		Gross Government Transportation (\(K^{\text{GT}}\)) & 4.07\% & -0.14\% & 4.44\% & 2.05\% & 1.06\% \\
		\bottomrule
	\end{tabular}
\end{table}

\paragraph{B}
These indicators show the structural breaks defining each regime. During the pre-war recovery (1931–1941), corporate GVA grew at a rapid 5.80\% per year. Public investment drove this, with infrastructure expanding at 4.07\% while private accumulation lagged. World War II changed everything. The mobilization block (1942–1945) redirected capital toward war production. Civilian infrastructure contracted at -0.14\% per year, while military spending pushed total GDP growth to a massive 9.92\%. The postwar Golden Age (1947–1973) represents the historical peak of coupled accumulation. Machinery grew at 8.49\% annually, structures at 6.38\%, and public infrastructure at 4.44\%. This balanced growth sustained a 4.98\% average expansion of real corporate GVA. In the post-Fordist era, this coordination broke down. Public infrastructure growth fell to 2.05\% before the GFC, and sank to a mere 1.06\% after 2008. Meanwhile, private structures decoupled, growing at 7.02\% pre-GFC before dropping to 3.99\% post-GFC. This combination of private overaccumulation and public decay explains the long-run slowdown of corporate capacity formation.
